\chapter*{پیشگفتار}
\addcontentsline{toc}{chapter}{پیشگفتار}
\thispagestyle{empty}
% ==========+==========+==========+==========+==========
 در حل مسایل عددی زبان های برنامه نویسی نقش تایین کننده ای در نحوه پیاده سازی الگوریتم ها و روش های عددی دارند. پارادایم زبان، توانایی تشخیص و مدیریت خطاهای عددی یا باگ ها، میزان استفاده از حافظه و پردازشگر تحت تاثیر مستقیم زبان مورد استفاده هستند.
اکثر زبان های مورد استفاده در محاسبات عددی از نوع 
\eng{imperative}
 هستند. در این میان زبان های 
\eng{functional}
 مثل 
\eng{Haskell}
 کمتر استفاده میشوند. در این گزارش ابتدا مقدمه ای بر این زبان ارایه میکنیم و سپس ویژگی هایی از آن را بررسی میکنیم که میتوانند در حل مسایل کاربردی باشند و در کنار
زبان هایی مانند 
\eng{MATLAB}
و
\eng{C}
 ابزار خوبی برای محققین باشد.
%===================================================== new line
زبان
\eng{Haskell}
ویژگی های منحصر به فردی دارد که در آنرا تبدیل به ابزار مناسبی برای پژوهش میکند.  
در
\cite{RefAHACwMG}
نویسندگان یک زبان
\eng{Domain Specific (DSL)}
با
\eng{Haskell}
میسازند تا یک زبان انتزاعی تر و ساده تر نسبت 
\eng{CUDA C++}
به دست آید. در نهایت با بررسی این زبان و مقایسه عملکرد آن با کد 
\eng{CUDA C++}
نتیجه میگرند که با کد های بسیار ساده تر میتوان به عملگردی مشابه به کد های 
\eng{CUDA C++}
دست یافت.
\eng{J. Karczmarczuk}
چند مقاله در زمینه مشتق گیری خودکار یا
\eng{Automatic Differentiation (AD)}
دارد که در آنها از زبان
\eng{Haskell}
برای پیاده سازی روش های خود استفاده کرده.  در 
\cite{RefOPFGP}
او از
\eng{Lazy Evaulation}
و تکنیک های آن در تولید و استفاده از دنباله های طویل یا
مشتقات عبارات عددی یک متغیره استفاده میکند.

